\scanpage{178}
% archive foliation: 177
\deleted{\unsure{Proposition 3}}

$\underline{\underline{\operatorname{Pic}}}_{X/S}=$ enveloppe exacte de $\underline{\underline{\mathrm{Prépic}}}_{X/S}$ (ou $\underline{\underline{\mathrm{Prépic}}}'_{X/S}$) pour la descente fid. plate quasi-compacte

[Il y a un ennui, car un \uline{recouvrement} ouvert d'un préschéma non quasi-cpt. n'est en général pas majoré par un morphisme fid. plat qu. cpt. Il faut donc dire au lieu de fid. plat qu. cpt: \deleted{fid plat séparé} \deleted{localisation} \uline{morphismes de localisation fidèlement plats}]

Donc on a des morphismes fonctoriels
\diagram{\begin{tikzcd}
\underline{\underline{\mathrm{Prépic}}}_{X/S}\arrow[r,"\alpha"]&\underline{\underline{\mathrm{Prépic}}}'_{X/S}\arrow[r,"\beta"]&\underline{\underline{\operatorname{Pic}}}_{X/S}
\end{tikzcd}}
Le premier morphisme n'est pas injectif en général [son noyau, si $\underline H^0_{X/S}(\underline O_X)=\underline O_S$, et si cela reste vrai par changement de base] est le foncteur $\underline{\underline{\operatorname{Pic}}}$ (absolu) en vertu de prop. 1, (i), \inserted{Il est surjectif} si on suppose en plus de la condition de [ ] que $X/S$ a une section. \deleted{$\unsure{\beta}$ \illegible{3}}
\margin{(\unsure{qui est non représentable par un schéma})}

